\documentclass[12pt]{article}
\usepackage[utf8]{inputenc}
\usepackage[T1]{fontenc}
\usepackage{hyperref}
\title{Mazaltepec Zapotec Language Resources for K-12 Curriculum Development}
\author{Compiled by Claude Sonnet 5.5 (claude-sonnet-5-5) for Andr\'as Kornai}
\date{}
\begin{document}
\maketitle

\section{Introduction}
Mazaltepec Zapotec (ISO 639-3 zpy, Glottocode maza1294) is a divergent Zapotec language of Oaxaca, Mexico. It is also called Etla Zapotec, Quiegolani Zapotec, Santo Tomas Mazaltepec and Western Yautepec Zapotec. Speaker estimates range from about 2,200 (Wikipedia, 1990 census) to 3,400 (Joshua Project). Wikipedia reports very low intelligibility with other Zapotec varieties. Little material exists beyond reference entries and religious audio.

\section{Primary Educational Resources}
No primers, textbooks or school programmes for this language were found in search results.

\section{Linguistic and Reference Materials}
\begin{itemize}
\item Wikipedia overview of the language: \url{https://en.wikipedia.org/wiki/Mazaltepec_Zapotec}
\end{itemize}

\section{Religious and Community Resources}
\begin{itemize}
\item Joshua Project profile, which reports that Bible portions were published in 2000 and that no New Testament or complete Bible exists: \url{https://joshuaproject.net/people_groups/13609}
\end{itemize}

\section{Digital and Audio Resources}
\begin{itemize}
\item Global Recordings Network, Zapotec, Mazaltepec: evangelism and Bible-teaching audio: \url{https://www.globalrecordings.net/en/language/zpy}
\end{itemize}

\section{Curriculum Development Context}
The language is reported to be used by all adults but not by all young people, and it is not known to be taught in schools. No dictionary, grammar or classroom material was found. A curriculum would need new documentation and materials development with the community. Because the variety is so different from other Zapotec languages, materials from other Zapotec varieties cannot simply be reused.

\end{document}